% !TeX root = ../main.tex
\section{Concluding remarks}
\label{sec:concluding-remarks}

In classical physics the potentials can be regarded as auxiliary objects,
since the force on a charge depends only on \(E\) and \(B\). In quantum
mechanics this is no longer the case. The Schrödinger equation of a charged
particle is written in terms of \(V\) and \(A\) through minimal coupling, and
the fields never appear in it explicitly. Gauge freedom survives, but only as
the joint freedom of redefining the potentials and the local phase convention
of the wavefunction.

What the potentials make observable is not their value at a point, which is
gauge-dependent, but the phase factor
\[
	\exp \left[ \frac{iq}{\hbar} \oint_C A \cdot dr \right]
	= \exp \left[ \frac{iq\Phi_B}{\hbar} \right],
\]
fixed by the flux linked by the loop \(C\). When \(B = 0\) throughout the
region accessible to the particle, this factor is unity in a simply connected
region, where the vector potential can be removed by a gauge transformation;
around a hole it can differ from unity, depending on the enclosed flux.
Wu~and~Yang~\cite{wuConceptNonintegrablePhase1975} later turned this observation into a general principle: the field strength
alone underdescribes electromagnetism, and the complete gauge-invariant
description is given by the phase factors around closed loops.

The prediction has been confirmed experimentally. After the early observations
of Chambers \cite{chambersShiftElectronInterference1960} and of M\"ollenstedt
and Bayh \cite{mllenstedtMessungKontinuierlichenPhasenschiebung1962}, the toroidal experiment of Tonomura
and collaborators removed the stray-field objection by confining the flux with a
superconducting layer, and observed relative phases of \(0\) or \(\pi\), as
required by flux quantization \cite{tonomuraEvidenceAharonovBohmEffect1986}.

The phase discussed here is also the first example of a broader class. Berry
\cite{berryQuantalPhaseFactors1984} showed that a quantum system transported adiabatically around a closed loop in
parameter space acquires a geometric phase, and treated the Aharonov--Bohm
phase as a particular case. Aharonov and
Casher \cite{aharonovTopologicalQuantumEffects1984} showed that a neutral particle with a magnetic moment, moving around a
line of charge, acquires an analogous topological phase, the dual of the
magnetic arrangement discussed above. Comprehensive reviews of the
theory and of the experiments can be found in
\cite{olariuQuantumEffectsElectromagnetic1985,peshkinAharonovBohmEffect1989}.
